% !TEX root = ../algebra_02.tex

\section{Теорема Гамильтона-Кэли}

Вычислим характеристический многочлен оператора, ограниченного на циклическое подпространство $C_v$. Матрица оператора $A|_{C_v}$ в базисе $v, Av, \dots, A^{d-1}v$
является сопровождающей матрицей для многочлена $M_{v,A}$:
\[
[A|_{C_v}] =
\begin{pmatrix}
	0 & 0 & \dots & 0 & -\beta_0 \\
	1 & 0 & \dots & 0 & -\beta_1 \\
	0 & 1 & \dots & 0 & -\beta_2 \\
	\vdots & \vdots & \ddots & \vdots & \vdots \\
	0 & 0 & \dots & 1 & -\beta_{d-1}
\end{pmatrix}.
\]

Будем использовать стандартное соглашение $\chi_A(x) = \det(xE - A).$
Тогда
\[
\chi_{A|_{C_v}}(x)
=
\det(xE - [A|_{C_v}])
=
x^d + \beta_{d-1}x^{d-1} + \dots + \beta_0
=
M_{v,A}(x).
\]

\begin{Theorem}{Гамильтона--Кэли}{}
	Пусть $\dim V < \infty$ и $A \in \operatorname{End} V$. Характеристический многочлен оператора $\chi_A(x)$ аннулирует этот оператор:
	\[
	\chi_A(A) = \mathbf{0}
	\].
\end{Theorem}

\begin{proof}
	Рассмотрим произвольный вектор $v \in V$ и порожденное им циклическое подпространство $C_v$.

	Характеристический многочлен оператора, ограниченного на инвариантное подпространство, делит характеристический многочлен всего оператора. Значит, $\chi_{A|_{C_v}}$ делит $\chi_A$.
	Поскольку матрица оператора на $C_v$ является сопровождающей матрицей минимального аннулятора, $\chi_{A|_{C_v}}$ совпадает с минимальным аннулятором вектора $v$, то есть $M_{v,A}$. Отсюда $\chi_A = M_{v,A} \cdot g$ для некоторого многочлена $g \in K[x]$.

	Подставим оператор $A$ и применим к вектору $v$:
	\[
	\chi_A(A)v = g(A)M_{v,A}(A)v = 0
	\],
	поскольку многочлен $M_{v,A}$ по определению аннулирует вектор $v$.

	Так как равенство $\chi_A(A)v = 0$ выполнено для абсолютно любого вектора $v \in V$, сам оператор $\chi_A(A)$ является нулевым: $\chi_A(A) = \mathbf{0}$.
\end{proof}
